\documentclass[11pt,a4paper]{article}
\usepackage[T1]{fontenc}
\usepackage[utf8]{inputenc}
\usepackage[czech]{babel}
\usepackage{lmodern}
\usepackage{microtype}
\usepackage{geometry}
\geometry{top=21mm,bottom=22mm,left=20mm,right=20mm}
\usepackage{amsmath,amssymb,mathtools}
\usepackage{enumitem,booktabs,multicol}
\usepackage{xcolor}
\usepackage[most]{tcolorbox}
\usepackage{tikz}
\usetikzlibrary{arrows.meta,positioning,calc,decorations.pathmorphing}
\usepackage{titlesec,fancyhdr,lastpage,hyperref}
\hypersetup{colorlinks=true,linkcolor=blue!50!black,urlcolor=blue!50!black}
\definecolor{VSNavy}{HTML}{17324D}
\definecolor{VSBlue}{HTML}{2F6B9A}
\definecolor{VSTeal}{HTML}{2A8C82}
\definecolor{VSGrey}{HTML}{576879}
\definecolor{VSLine}{HTML}{DBE3EB}
\definecolor{VSFog}{HTML}{F4F7FA}
\setlength{\parindent}{0pt}
\setlength{\parskip}{5.2pt}
\setlist[itemize]{topsep=3pt,itemsep=2pt,leftmargin=1.3em}
\setlist[enumerate]{topsep=4pt,itemsep=6pt,leftmargin=1.55em}
\titleformat{\section}{\Large\bfseries\color{VSNavy}}{\thesection.}{.75em}{}
\titleformat{\subsection}{\large\bfseries\color{VSBlue}}{\thesubsection}{.7em}{}
\titlespacing*{\section}{0pt}{16pt}{8pt}
\titlespacing*{\subsection}{0pt}{12pt}{6pt}
\pagestyle{fancy}\fancyhf{}
\fancyhead[L]{\footnotesize\color{VSNavy}\textbf{VOJTĚCH STEIDL}}
\fancyhead[R]{\footnotesize\color{VSGrey}Fyzika \,/\, termika}
\fancyfoot[L]{\footnotesize\color{VSGrey}Výukový zápis a řešené příklady}
\fancyfoot[R]{\footnotesize\color{VSGrey}\thepage\,/\,\pageref{LastPage}}
\renewcommand{\headrulewidth}{.35pt}
\renewcommand{\footrulewidth}{0pt}
\newtcolorbox{theorybox}[1][]{enhanced,breakable,colback=VSFog,colframe=VSLine,boxrule=.55pt,arc=2mm,left=3mm,right=3mm,top=2mm,bottom=2mm,#1}
\newtcolorbox{solvedbox}[1][]{enhanced,breakable,colback=white,colframe=VSTeal!70!black,boxrule=.8pt,arc=1.5mm,left=3.5mm,right=3.5mm,top=2.5mm,bottom=2.5mm,#1}
\newif\ifshowanswers
\showanswerstrue % Pro verzi bez závěrečných výsledků změň na \showanswersfalse.
\newcommand{\unit}[1]{\,\mathrm{#1}}
\begin{document}

{\color{VSBlue}\footnotesize\bfseries FYZIKA SŠ \quad / \quad TERMODYNAMIKA \quad / \quad STUDIJNÍ ZÁPIS}\par
\vspace{4mm}
{\Huge\bfseries\color{VSNavy}Vnitřní energie a přenos tepla}\par
{\Large\color{VSBlue}Vedení, proudění, záření a teplotní roztažnost}\par
\vspace{3mm}
{\color{VSGrey}\rule{\linewidth}{.6pt}}\par
\vspace{2mm}
\textbf{Co budeš po lekci umět:} rozlišit teplotu, vnitřní energii, teplo a práci; určit směr přenosu energie; rozpoznat vedení, proudění a záření; řešit úlohy na $Q=mc\Delta T$ a délkovou roztažnost.

\begin{theorybox}
\textbf{Čtyři pojmy, které se nesmějí plést}

\begin{tabular}{@{}p{.18\linewidth}p{.76\linewidth}@{}}
\toprule
\textbf{Pojem} & \textbf{Přesný význam} \\
\midrule
Teplota $T$ & Stavová veličina vyjadřující tepelný stav tělesa; v kelvinech K, běžně také v \,$^\circ$C.\\
Vnitřní energie $U$ & Energie související s mikroskopickým pohybem a vzájemným působením částic tělesa. Jednotka J.\\
Teplo $Q$ & Energie přenesená \emph{v důsledku teplotního rozdílu}. Není to zásoba „uvnitř“ tělesa. Jednotka J.\\
Práce $W$ & Jiný způsob předávání energie, např. stlačováním plynu nebo třením. Jednotka J.\\
\bottomrule
\end{tabular}
\end{theorybox}

\section{Vnitřní energie a její změna}
Vnitřní energie zahrnuje mikroskopickou kinetickou a potenciální energii částic. Může se změnit \textbf{konáním práce}, \textbf{přenosem tepla} nebo oběma způsoby současně.

Pro zvolenou znaménkovou konvenci, kdy $W_{\mathrm{na}}$ značí \emph{práci vykonanou na soustavě}, platí:
\[
\boxed{\Delta U=Q+W_{\mathrm{na}}}
\]
Kladné $Q$ znamená teplo přijaté soustavou; záporné $Q$ teplo odevzdané. Při stlačování plynu okolím je $W_{\mathrm{na}}>0$.
Některé učebnice užívají práci vykonanou \emph{soustavou} $W_{\mathrm{od}}$ a zapisují $\Delta U=Q-W_{\mathrm{od}}$. Oba zápisy jsou ekvivalentní.

\textbf{Příklady změny vnitřní energie:}
\begin{itemize}
\item Třením dlaní konáme práci; dlaně se ohřejí.
\item Kovovou lžíci ponoříme do teplé vody; energii získává přenosem tepla.
\item Rychlým stlačením vzduchu v pumpičce konáme práci; vzduch se může zahřát.
\end{itemize}

\subsection{Kolik tepla je potřeba k ohřevu?}
\begin{theorybox}
Při změně teploty \textbf{bez změny skupenství} a za předpokladu přibližně konstantní měrné tepelné kapacity:
\[
\boxed{Q=mc\Delta T},\qquad \Delta T=T_2-T_1.
\]
$m$ je hmotnost v kg; $c$ měrná tepelná kapacita v $\mathrm{J\,kg^{-1}\,K^{-1}}$; rozdíl teplot $\Delta T$ lze vyjádřit ve stupních Celsia i v kelvinech (číselně se shoduje).
Pro vodu v běžném školním výpočtu bereme $c_{\mathrm{voda}}\approx 4200\unit{J\,kg^{-1}\,K^{-1}}$.
\end{theorybox}

\begin{solvedbox}
\textbf{Řešený příklad 1: Ohřev vody}\\
Jakou energii potřebujeme pro ohřátí $0{,}25\unit{kg}$ vody z $20\,^\circ\mathrm C$ na $60\,^\circ\mathrm C$? Zanedbej tepelné ztráty.

\textbf{1. Známé veličiny:} $m=0{,}25\unit{kg}$, $c=4200\unit{J\,kg^{-1}K^{-1}}$ a $\Delta T=60-20=40\unit{K}$.

\textbf{2. Dosazení s jednotkami:}
\begin{align*}
Q&=mc\Delta T\\
 &=0{,}25\unit{kg}\cdot4200\unit{J\,kg^{-1}K^{-1}}\cdot40\unit{K}\\
 &=42\,000\unit{J}.
\end{align*}
\textbf{3. Výsledek:}\quad $\boxed{Q=42\unit{kJ}}$.
\end{solvedbox}

\section{Tři způsoby přenosu tepla}
Energie se při tepelném kontaktu samovolně přenáší z oblasti \textbf{vyšší teploty} do oblasti \textbf{nižší teploty}. Rozdíl je v \emph{mechanismu} přenosu.

\subsection{Vedení (kondukce)}
Při \textbf{vedení} se energie předává mikroskopickými interakcemi mezi sousedními částicemi \emph{bez makroskopického proudění látky}. U kovů se významně uplatňují i volné elektrony.

\begin{center}
\begin{tikzpicture}[>=Latex,font=\small]
\fill[orange!45] (0,0) rectangle (1.4,.72);
\fill[blue!17] (1.4,0) rectangle (7.0,.72);
\draw[VSNavy] (0,0) rectangle (7,.72);
\foreach \x in {2,3,4,5,6} {\draw[-{Latex[length=2mm]},very thick,VSTeal] (\x-.25,.36)--(\x+.38,.36);}
\node[below] at (.7,-.12) {teplý konec};
\node[below] at (6,-.12) {chladný konec};
\node[above] at (3.5,.78) {kovová tyč: přenos energie};
\end{tikzpicture}
\end{center}
\textbf{Příklady:} zahřívání rukojeti kovové lžíce; tepelný most v konstrukci domu. Dřevo a vzduch bývají horší vodiče než kovy.

Pro rovinnou vrstvu při ustáleném jednorozměrném vedení přibližně platí:
\[
\boxed{P=\frac{Q}{t}=\lambda A\frac{T_{\mathrm{teplá}}-T_{\mathrm{studená}}}{d}}
\]
kde $\lambda$ je tepelná vodivost v $\mathrm{W\,m^{-1}\,K^{-1}}$, $A$ plocha a $d$ tloušťka vrstvy. \textbf{Silnější izolace} za stejných podmínek snižuje tepelný tok.

\subsection{Proudění (konvekce)}
Při \textbf{proudění} přenáší energii pohybující se kapalina nebo plyn. Při přirozené konvekci vzniká pohyb vlivem rozdílů hustoty a vztlaku; u nucené konvekce jej vytváří čerpadlo či ventilátor.

\begin{center}
\begin{tikzpicture}[>=Latex,font=\small]
\draw[VSNavy,thick,rounded corners=2pt] (0,0) rectangle (7.1,2.05);
\fill[orange!25] (.15,.12) rectangle (6.95,.35);
\draw[->,very thick,red!75!black] (1.4,.45) .. controls (1.05,1.35) and (1.65,1.55) .. (2.7,1.7);
\draw[->,very thick,red!75!black] (3.2,.5) .. controls (3.35,1.4) and (4,1.55) .. (4.5,1.68);
\draw[->,very thick,blue!70!black] (5.7,1.72) .. controls (6,1.25) and (5.95,.8) .. (5.4,.46);
\draw[->,very thick,blue!70!black] (.65,1.7) .. controls (.45,1.2) and (.55,.7) .. (1,.42);
\node[below] at (3.55,-.13) {ohřev u dna nádoby};
\node[above] at (3.55,2.07) {proudění kapaliny};
\end{tikzpicture}
\end{center}
\textbf{Příklady:} cirkulace vody v hrnci; stoupající vzduch nad radiátorem; proudění vzduchu v místnosti. \textbf{Konvekce neprobíhá ve vakuu}, protože vyžaduje látkové prostředí.

\subsection{Záření (radiace)}
Každé těleso s teplotou vyšší než $0\unit K$ vyzařuje elektromagnetické záření. K přenosu \textbf{není potřeba látkové prostředí}: energie ze Slunce dopadá na Zemi i přes vakuum.

\textbf{Příklady:} sluneční záření; teplo vnímané u rozpálených kamen; infračervená kamera. Tmavý a matný povrch zpravidla více pohlcuje dopadající záření než lesklý světlý povrch (při porovnatelné vlnové délce a materiálu).

\begin{theorybox}
Pro ideální černé těleso platí \textbf{Stefanův--Boltzmannův zákon}:
\[
P=\sigma A T^4.
\]
U reálného šedého povrchu používáme $P=\varepsilon\sigma A T^4$, kde $0\leq\varepsilon\leq1$ je emisivita. Teplota musí být v \textbf{kelvinech}. Pro čistý výsledný radiační tok s okolím o teplotě $T_{\mathrm{okolí}}$ přibližně platí $P_{\mathrm{net}}=\varepsilon\sigma A(T^4-T_{\mathrm{okolí}}^4)$.
\end{theorybox}
\textbf{Pozor:} Záření, vedení a proudění mohou probíhat \emph{současně} (např. radiátor předává energii vzduchu prouděním a okolním tělesům zářením).

\begin{solvedbox}
\textbf{Řešený příklad 2: Zateplení stěny}\\
Izolace má vodivost $\lambda=0{,}04\unit{W\,m^{-1}K^{-1}}$, plochu $A=1\unit{m^2}$, tloušťku $d=0{,}10\unit m$ a rozdíl teplot mezi povrchy $\Delta T=20\unit K$. Urči výkon přenášený vedením. Předpokládej ustálené vedení a zanedbej povrchové odpory.
\begin{align*}
P&=\lambda A\frac{\Delta T}{d}\\
 &=0{,}04\frac{\mathrm W}{\mathrm{m\,K}}\cdot1\unit{m^2}\cdot
   \frac{20\unit K}{0{,}10\unit m}\\
 &=\boxed{8\unit W}.
\end{align*}
Při \emph{dvojnásobné} tloušťce izolace a stejných povrchových teplotách by byl výsledek $4\unit W$.
\end{solvedbox}

\section{Teplotní roztažnost}
Při zahřívání se rozměry většiny pevných látek zvětšují. Zahřátý materiál se v průměru roztahuje, protože se mění rovnovážné vzdálenosti částic.

\begin{theorybox}
Pro malou změnu teploty a přibližně konstantní koeficient délkové roztažnosti $\alpha$:
\[
\boxed{\Delta l=\alpha\, l_0\,\Delta T},\qquad
\boxed{l=l_0(1+\alpha\Delta T)}.
\]
$\alpha$ má jednotku $\mathrm{K^{-1}}$; $l_0$ je původní délka. Pokud $\Delta T<0$, jde o \textbf{smrštění}.
\end{theorybox}
Příklady: dilatační spáry mostů a kolejí, změna délky kovových trubek, bimetalický pásek termostatu.

U izotropních pevných látek pro malé teplotní změny přibližně platí plošná roztažnost $\Delta S\approx 2\alpha S_0\Delta T$ a objemová $\Delta V\approx 3\alpha V_0\Delta T$.
U kapalin se zpravidla sleduje roztažnost objemová. \textbf{Voda mezi $0$ a $4\,^\circ$C} představuje důležitou výjimku: při ohřevu směrem ke $4\,^\circ$C její hustota roste.

\begin{solvedbox}
\textbf{Řešený příklad 3: Ocelová tyč}\\
Ocelová tyč má délku $l_0=2{,}0\unit m$. Při ohřevu o $50\unit K$ urči prodloužení, jestliže $\alpha=12\cdot10^{-6}\unit{K^{-1}}$.
\begin{align*}
\Delta l&=\alpha l_0\Delta T\\
 &=12\cdot10^{-6}\unit{K^{-1}}\cdot2{,}0\unit m\cdot50\unit K\\
 &=0{,}0012\unit m=\boxed{1{,}2\unit{mm}}.
\end{align*}
Nová délka je $l=\boxed{2{,}0012\unit m}$.
\end{solvedbox}

\section{Procvičení bez nápovědy}
\textit{Nejdřív zkus určit mechanismus nebo vztah, teprve potom počítej.}
\begin{enumerate}
\item Kovovou lžíci necháme v horké polévce. Jaký přenos energie je dominantní uvnitř kovové rukojeti?
\item Proč se nad zapnutým radiátorem často zvedá teplejší vzduch? Popiš, proč jde o proudění.
\item Jak se ze Slunce přenese energie na Zemi? Proč to nemůže být vedení ani proudění ve vakuu?
\item Vypočti teplo potřebné k ohřevu $0{,}5\unit{kg}$ vody z $15\,^\circ$C na $35\,^\circ$C. Uvažuj $c=4200\unit{J\,kg^{-1}\,K^{-1}}$.
\item Ustáleným vedením prochází vrstvou při nezměněných podmínkách výkon $20\unit W$. Jaký bude při zdvojnásobení tloušťky?
\item Tyč o délce $1{,}5\unit m$ má $\alpha=20\cdot10^{-6}\unit{K^{-1}}$. O kolik se prodlouží při ohřevu o $40\unit K$?
\item Povrch tělesa má dvakrát vyšší absolutní teplotu než dříve. Kolikrát větší je jeho \emph{celkový vyzařovaný výkon}, pokud se nemění emisivita a plocha?
\item Vysvětli, proč není přesné tvrzení „těleso obsahuje teplo“.
\end{enumerate}

\ifshowanswers
\section{Výsledky a odůvodnění}
\begin{enumerate}
\item \textbf{Vedení.} Energie postupuje materiálem bez proudění kovu jako celku.
\item \textbf{Přirozená konvekce.} Zahřátý vzduch má při přibližně stejném tlaku menší hustotu a působením vztlaku stoupá; chladnější jej nahrazuje.
\item \textbf{Elektromagnetickým zářením.} Vakuum neobsahuje látkové prostředí potřebné pro vedení nebo proudění.
\item $\Delta T=20\unit K$;
\[
Q=0{,}5\cdot4200\cdot20=\boxed{42\,000\unit J=42\unit{kJ}}.
\]
\item \textbf{$10\unit W$.} V jednoduchém modelu je $P$ nepřímo úměrný tloušťce $d$.
\item
\[
\Delta l=20\cdot10^{-6}\cdot1{,}5\cdot40
 =0{,}0012\unit m=\boxed{1{,}2\unit{mm}}.
\]
\item \textbf{$2^4=16$krát.} Ve Stefanově--Boltzmannově zákoně je $T^4$, nikoli $T^2$. Nelze ale bez dalšího tvrdit, že stejně roste \emph{čistý} tok energie do okolí.
\item \textbf{Teplo} označuje energii \emph{přenášenou} mezi soustavami kvůli rozdílu teplot; energetický stav soustavy popisuje mimo jiné její \textbf{vnitřní energie} $U$.
\end{enumerate}
\fi

\begin{theorybox}
\textbf{Shrnutí v jedné minutě}\\
\textbf{Vedení:} energie přes sousední částice. \quad
\textbf{Proudění:} s pohybem kapaliny nebo plynu.\\
\textbf{Záření:} elektromagnetické vlny, funguje i ve vakuu.\\
\textbf{Výpočty:} $Q=mc\Delta T$, $P=\lambda A\Delta T/d$ (ustálené vedení),
$\Delta l=\alpha l_0\Delta T$.\\
\textbf{Kontrola:} teplotní rozdíl je v K nebo $^\circ$C stejný, ale ve vzorci pro záření používej absolutní teplotu v \emph{kelvinech}.
\end{theorybox}
\textbf{Další krok:} interaktivní laboratoř a test \href{https://vojtechsteidl.eu/interactive-notes/thermal-energy/}{\texttt{vojtechsteidl.eu/interactive-notes/thermal-energy/}}.
\end{document}
